% RESEARCH_DRAFT_SHA256: f1b8b51f2c970399b6d488a71f17cf649be03934239d63090077a9c2e5e23523
% Layout and source-reference wrappers expanded for the web renderer.
\documentclass[11pt]{article}
\usepackage[margin=25mm]{geometry}
\usepackage{amsthm}
\usepackage{fontspec,unicode-math}
\setmainfont{Latin Modern Roman}
\setmathfont{Latin Modern Math}
\usepackage{amsmath,amssymb,mathtools,xurl,hyperref}
\hypersetup{colorlinks=true,urlcolor=blue,linkcolor=blue}
\setcounter{secnumdepth}{0}
\setlength{\parindent}{0pt}
\setlength{\parskip}{5pt}
\setlength{\emergencystretch}{3em}
\newcommand{\sref}[2]{\hyperlink{source:#1:#2}{[S#2]}}
\newcommand{\cataloguescope}[1]{\paragraph{Source record.}#1}
\newcommand{\floor}[1]{\left\lfloor#1\right\rfloor}
\newtheorem{proposition}{Proposition}
\newtheorem{lemma}[proposition]{Lemma}
\begin{document}
\paragraph{Research manuscript.} OpenAI \texttt{gpt-6-astra}, reasoning \texttt{ultra}; 11 October 2026.
\paragraph{Preserved source body.} The complete catalogue mathematical text and sources follow with layout and source links expanded, including source qualifications and its historical preparation notes. New research begins at the section headed Research attempt.

% UnsolvedMath ID 30004101; source catalogue code OWR-16775-003
\section{Sequential Decoding of Nonbinary Convolutional Codes}\label{30004101}
{\small UnsolvedMath ID 30004101 · Computer Science}
\subsection{Definitions and mathematical statement}

\paragraph{Shared notation.}Standard number systems are $\mathbb N,\mathbb Z,\mathbb Q,\mathbb R,\mathbb C$. Specialized notation and index ranges follow the cited source and the definitions below.


A nonbinary convolutional code encodes a stream of symbols over a finite field $\mathbb F_q$, $q>2$, by a finite-memory linear encoder. Its trellis records encoder states and possible output transitions. A sequential decoder explores paths in this trellis using a metric derived from the noisy channel. Relevant performance measures include decoding error probability and the distribution of decoding work at a prescribed rate.
\paragraph{Record scope and qualifications.}
The catalogue title is retained as metadata. Named certificates, contracts, bounds or frameworks in that title are not asserted as results of the original source extract.
This is a family of performance questions: the report does not specify one channel, encoder, decoding metric, or error criterion. Any quantitative answer must state those choices.

Source-record limitation: The nonbinary convolutional-code family, channel/noise regime, sequential-decoding method, and performance target are not supplied.
\paragraph{Statement recorded in the supplied source.}
How do nonbinary convolutional codes perform with sequential decoders on channels with errors? \href{https://doi.org/10.4171/owr/2019/13}{[S2]}
\paragraph{Source status.}
The supplied catalogue labels this source record open; that metadata was checked against the supplied snapshot on 2026-10-11. The label is preserved for provenance and is not an independent verification of present open status or mathematical correctness. The catalogue literature assessment is dated 2026-08-22.
\subsection{Short English statement}
Analyze reliability and decoding effort for sequential decoding of convolutional codes over larger alphabets.
\subsection{Sources}
\begin{itemize}
\item[S1] Ulam AI and the UnsolvedMath Contributors, UnsolvedMath: A Curated Collection of Open Mathematics Problems, record 30004101 (OWR-16775-003). Catalogue curation and metadata: CC BY 4.0. The source extract is retained; definitions, English summaries and qualifications are editorial additions.
\url{https://huggingface.co/datasets/ulamai/UnsolvedMath}.
\item[S2] Contemporary Coding Theory (2020), Oberwolfach Reports, DOI 10.4171/owr/2019/13, pp. 773--840
\url{https://doi.org/10.4171/owr/2019/13}.
\item[S3] Sequential decoding for synchronization errors
\url{https://arxiv.org/abs/2201.11935}.
\item[S4] Multi-sequence sequential decoding
\url{https://doi.org/10.1109/TCOMM.2024.3405322}.
\item Source-specific Oberwolfach Reports licences, including CC BY-SA 4.0 where applicable, are retained; see THIRD\_PARTY\_NOTICES.md.
\end{itemize}
\paragraph{Preparation and review.}
The original source and UnsolvedMath contributors retain attribution for the source extract. Codex (GPT-6; exact serving revision not exposed) assisted with definitions, English summaries, scope notes and formatting. No new human mathematical authorship or independently verified solution is asserted. Permanent source ID 30004101 and code OWR-16775-003 checked against the complete supplied registry and current repository. Distinct upstream records, including repeated questions and broad or damaged extracts, are retained for complete coverage of the source CS category. A repeated record is not a new independent conjecture. Complete mathematical-equivalence and current-literature review has not been performed. Preparation and final review follow contribution prompts C and D, with source gaps disclosed. Structural and rendering checks do not establish a new theorem or open status.
\label{end:30004101}

\section{Research attempt}
\subsection{Literature, source context and scope}
The full original catalogue body and source qualifications are preserved above. The relevant primary workshop passages were read in the cached publisher report. Historical claims and subsequent attributed results are distinguished from the mathematical work below. A restricted example, standard lemma or known theorem is not counted as new progress on the entire original target. Missing assumptions and repeated or rhetorical extracts are not silently turned into different solved problems.

\subsection{A terminated trellis baseline isolates the harder search question}
The source distinguishes sequential tree-search decoders from Viterbi decoding and emphasizes dependence on the column-distance profile. We analyze a small fixed-memory baseline without calling it a general sequential-decoding theorem. Over $\mathbb F_3$, take $T\ge1$ information symbols $u_1,\ldots,u_T$, put $u_0=u_{T+1}=0$, and transmit the $T+1$ pairs
\[
c_t=(u_t,u_t+u_{t-1}),\qquad 1\le t\le T+1.
\]
Assume independent ternary symmetric substitution noise: a symbol is unchanged with probability $1-p$ and becomes each other symbol with probability $p/2$, where $0<p<2/3$.
\begin{proposition}
This terminated code has minimum Hamming distance 3. Maximum-likelihood decoding can be performed by a three-state trellis in $O(T)$ arithmetic operations, and it corrects every single-symbol error uniquely.
\end{proposition}
\begin{proof}
For the difference of two input sequences, let $r$ and $s$ be the first and last nonzero positions. The systematic coordinates contain at least one nonzero symbol. The parity coordinate at $r$ is nonzero because its preceding difference is zero, and that at $s+1$ is nonzero because its following difference is zero. These are distinct coordinates, so the output weight is at least 3. A single nonzero input symbol attains weight 3, proving equality. Two different codewords cannot both be within distance one of a received word, by the triangle inequality, proving unique correction.
A received length-$2(T+1)$ word has likelihood $(1-p)^{L-h}(p/2)^h$ at Hamming distance $h$ from a codeword, where $L=2(T+1)$. Since $1-p>p/2$, maximizing likelihood is minimizing Hamming distance. Let $D_t(s)$ be the smallest accumulated distance among paths ending in state $u_t=s$. Initialize $D_0(0)=0$ and the other states to infinity, and update
\[
D_t(s)=\min_{r\in\mathbb F_3}\{D_{t-1}(r)+d_H(y_t,(s,s+r))\}.
\]
Induction on $t$ proves this minimum-path interpretation: every path has a unique previous state, and the displayed transition is its exact new distance cost. At the final stage require $s=0$. Back pointers recover a minimizing path. Each stage considers nine transitions, proving the runtime bound.
\end{proof}
\subsection{Verification and remaining gaps}
For $T=3$, all 27 codewords and all their single-symbol corruptions were enumerated, and the trellis recovered the original input uniquely. The minimum pairwise distance was checked to be 3. This efficient baseline uses only three states; it does not address growing encoder memory, the random amount of work of stack or Fano sequential search, cutoff-rate behavior, or a sharp error-probability tradeoff. The source explicitly warns that larger state spaces change the comparison. No general answer to its nonbinary sequential-decoding agenda is obtained.

\section{References and source attribution}
Primary workshop report: \url{https://doi.org/10.4171/owr/2019/13}. Relevant passages read: 2019/13 pp. 793--794, nonbinary convolutional codes, Viterbi versus sequential decoding and the column-distance question.

All retrievals were checked on 11 October 2026. Original question attribution remains in the preserved Sources section. External theorems are attributed and not claimed as new discoveries or fully verified merely by citation.


\paragraph{Integration review.} Codex (GPT-6; exact serving revision not exposed) reviewed the written arguments and expanded layout and source-reference wrappers. Repository metadata and the production web renderer were checked using Python 3.13 and JavaScript with MathJax 3.2.2. This model-assisted review is not independent human proof verification; no PDF compilation or formal proof checking is asserted.

\section{Completion Estimate}
\noindent\textbf{COMPLETION: 0\%}
No substantiated new progress toward the entire original target is claimed. The worked results and failed route are delimited above; a known result, illustrative model, or auxiliary obstruction is not counted as original-target completion.

\section{Final status and attribution}
\textbf{Attempt status: unresolved.} This describes this research attempt, not a new ruling on the source catalogue's status or on external literature claims.
The mathematics and this self-audit were generated by OpenAI \texttt{gpt-6-astra}, reasoning \texttt{ultra}, on 11 October 2026. Codex (GPT-6; exact serving revision not exposed) orchestrated the batch, reviewed the written mathematical arguments, and contributed formatting/repository review. This additional model review is not independent human verification. No human mathematical authorship, independent human proof checking, or formal verification is claimed. The historical preparation paragraphs in the reproduced catalogue body concern that body only.

Contribution prompt A was applied to the complete definition and supplied source record, followed by the adversarial checks of prompt B. Local reproduction files are \texttt{batch022/checks.py}, \texttt{computations.json}, and the per-ID audit JSON. They were not public links at preparation time. Numerical checks are finite and do not prove unrestricted results. No LaTeX engine was found locally; compilation was not performed by this research agent.

\end{document}
